\begin{figure*}[!t]
	\centering
	\includegraphics[width=7.0in]{Fig1_PIML_KD_Framework_V2-1_2.pdf}
	\caption{Architectural overview of the proposed \textit{MicroPhys-BMS} framework. (a) Real-world operational vehicle microtrip current profile and transient relaxation voltage dynamics ($30\text{--}600\text{ s}$). (b) Domain-guided physical feature engineering extracting 12-dimensional electro-thermal descriptors ($X_{\text{PIML}}$). (c) Hybrid knowledge distillation compressing a 300-tree LightGBM ensemble into a 961-parameter student Micro-MLP under multi-objective physics-calibrated loss. (d) Automotive embedded deployment pipeline yielding a MISRA-C compliant, zero-\texttt{malloc} static C header (\texttt{bms\_soc\_piml\_mlp3.h}) for deterministic execution on automotive microcontrollers.}
	\label{fig:framework_overview}
\end{figure*}

\section{Introduction}
\label{sec:introduction}

\IEEEPARstart{L}{ithium-ion} batteries, particularly Lithium Iron Phosphate ($\text{LiFePO}_4$ or LFP), have established decisive dominance as the preferred electrochemical energy storage technology for commercial electric vehicles (EVs), urban transit buses, and heavy-duty grid energy storage systems \cite{masih2013optimum, rahimirad2019battery, attia2020closed, rout2024online}.

\subsection{Industrial Context and the LFP Flat-Plateau Challenge}
\label{sec:intro_a}

Driven by their exceptional inherent thermal runaway resilience, operational cycle life exceeding 3000 full equivalent cycles, non-toxic cobalt- and nickel-free chemistry, and superior raw material cost-effectiveness, LFP battery packs have achieved dominant market penetration in commercial transportation electrification \cite{masih2013optimum, rahimirad2019battery, attia2020closed, severson2019data, najafi2025adaptive}. Despite these compelling safety and economic advantages, deploying LFP battery chemistries in advanced automotive battery management systems (BMS) introduces severe, non-linear state observation challenges that are fundamentally distinct from nickel-manganese-cobalt (NMC) or lithium titanate (LTO) chemistries \cite{roscher2011dynamic, hu2012comparative, sorouri2025breaking, wang2025lifepo4}.

The central impediment to reliable on-board state-of-charge (SOC) tracking in LFP cells resides in the thermodynamic flat open-circuit voltage (OCV) plateau, which spans an extensive operational corridor from approximately $20\%$ to $80\%$ SOC \cite{sorouri2025breaking, xiong2018towards}. From a solid-state electrochemical perspective, lithium deintercalation and insertion in the olivine $\text{LiFePO}_4$ lattice proceed via a first-order phase transformation between two distinct ordered crystal phases: the fully lithiated triphylite ($\text{LiFePO}_4$) phase and the fully delithiated heterosite ($\text{FePO}_4$) phase \cite{sorouri2025breaking}. In the thermodynamic miscibility gap, the chemical potential of lithium ions $\mu_{\text{Li}^+}$ remains strictly constant across the phase boundary. According to the Nernst-Planck thermodynamic relation:
\begin{equation}
	\label{eq:nernst_ocv}
	V_{\text{ocv}}(\text{SOC}, T_K) = -\frac{\mu_{\text{Li}^+}^{\text{cathode}}(\text{SOC}, T_K) - \mu_{\text{Li}^+}^{\text{anode}}(\text{SOC}, T_K)}{F},
\end{equation}
where $F = 96485.33\ \text{C/mol}$ denotes Faraday's constant and $T_K$ is absolute temperature. Consequently, as illustrated in Fig.~\ref{fig:physical_motivation}(a), the incremental derivative of the equilibrium open-circuit voltage with respect to SOC is virtually zero throughout this wide domain:
\begin{equation}
	\label{eq:flat_plateau_derivative}
	\left. \frac{\partial V_{\text{ocv}}(\text{SOC}, T_K)}{\partial \text{SOC}} \right|_{\text{SOC} \in [0.20, 0.80]} \approx 0.10 \text{--} 0.50\ \text{mV/\%}.
\end{equation}

\begin{figure}[!t]
	\centering
	\includegraphics[width=\columnwidth]{Figure_2_Physical_Motivation_Final_IEEE_2.pdf}
	\caption{Electrochemical characteristics and observability dilemma of LFP cells. (a) Static $V_{\text{ocv}}$--SOC characteristic across $10^\circ\text{C}$, $25^\circ\text{C}$, and $45^\circ\text{C}$ highlighting the flat plateau. (b) Dynamic terminal voltage relaxation trajectories across varied SOC levels ($20\%$, $40\%$, $60\%$, and $80\%$). (c) Direct thermodynamic inversion error under typical analog front-end (AFE) sensing offsets ($+2\text{ mV}$, $+5\text{ mV}$, and $+10\text{ mV}$).}
	\label{fig:physical_motivation}
\end{figure}

Compounding this flat derivative is significant path-dependent electrochemical hysteresis between charging and discharging branches, which arises from mechanical coherency strain energy and phase-boundary interfacial resistance \cite{roscher2011dynamic, sorouri2025breaking}. Because the static voltage curve is nearly horizontal, any direct numerical inversion of the measured terminal voltage $V_{\text{mea}}$ to deduce battery state exhibits extreme mathematical ill-conditioning. As demonstrated in Fig.~\ref{fig:physical_motivation}(c), an analog front-end (AFE) measurement offset of merely $+5.0\ \text{mV}$ induces catastrophic SOC estimation errors exceeding $20\%\text{--}30\%$ across the plateau domain.

Furthermore, dynamic vehicle operation introduces prolonged, non-linear voltage relaxation tails following current pulses (Fig.~\ref{fig:physical_motivation}(b)). Following current cessation during real-world vehicle microtrips (such as regenerative braking events or urban traffic stops), the terminal voltage does not instantaneously recover to thermodynamic equilibrium $V_{\text{ocv}}$. Instead, it undergoes a complex, multi-time-scale relaxation process governed by fast double-layer capacitive discharge ($\tau_1 \le 10\ \text{s}$) and sluggish solid-state lithium-ion diffusion within the porous active electrode particles ($\tau_2 \ge 100\ \text{s}$), lasting up to several hours \cite{che2025enhanced, sorouri2025breaking}.

In commercial automotive BMS applications, short rest intervals ($30\ \text{s}$ to $600\ \text{s}$) represent crucial operational opportunities for re-initializing and resetting the cumulative integration drift inherent to Coulomb counting algorithms \cite{waag2014critical, che2025enhanced}. However, sampling the terminal voltage prematurely during non-equilibrium relaxation windows without accurately back-extrapolating the true equilibrium overpotential introduces large state initialization errors. Developing a computationally deterministic and mathematically bounded algorithm capable of extracting equilibrium signatures from short-term relaxation telemetry is therefore vital for next-generation transportation electrification.

\subsection{Limitations of Classical Observers and Heavy Tree Ensembles}
\label{sec:intro_b}

To track battery state dynamics under operational driving loads, conventional automotive BMS architectures deploy state estimation algorithms falling into two broad paradigms: equivalent circuit model (ECM)-based adaptive observers and data-driven machine learning models \cite{masih2018novel, waag2014critical, najafi2024statepower, xu2025state}.

Classical state observation frameworks rely heavily on Extended Kalman Filters (EKF), Unscented Kalman Filters (UKF), Adaptive Extended Kalman Filters (AEKF), or recursive least squares (RLS) coupled with low-order RC equivalent circuits \cite{plett2004extended, xiong2018towards, wei2022future, hou2026online, fan2025state}. In these recursive formulations, the state correction step is strictly driven by the voltage innovation residual $\tilde{y}_k = V_{\text{mea},k} - \hat{V}_{t,k}$, weighted through the Kalman gain matrix $\mathbf{K}_k$ \cite{plett2004extended, wei2022future}:
Classical state observation frameworks rely on Extended Kalman Filters (EKF), Unscented Kalman Filters (UKF), Adaptive EKF, or recursive least squares (RLS) coupled with low-order RC equivalent circuits \cite{plett2004extended, xiong2018towards, wei2022future, hou2026online, fan2025state}. The state correction step uses the voltage innovation residual weighted through the Kalman gain $\mathbf{K}_k = \mathbf{P}_{k|k-1}\mathbf{H}_k^T(\mathbf{H}_k\mathbf{P}_{k|k-1}\mathbf{H}_k^T + \mathbf{R}_v)^{-1}$, where $\mathbf{H}_k = \partial V_{\text{ocv}}/\partial \text{SOC}$ is the observation Jacobian. Within the LFP flat plateau ($\text{SOC} \in [0.20, 0.80]$), $\mathbf{H}_k \to \mathbf{0}$, causing $\mathbf{K}_k \to \mathbf{0}$ \cite{wei2022future, sorouri2025breaking}. Feedback correction collapses into open-loop Coulomb counting $\text{SOC}_k = \text{SOC}_{k-1} - (\eta_c \Delta t / Q_{\text{nom}})(I_k + \delta I_k)$, where $\delta I_k$ captures sensor bias and ADC quantization errors; accumulated state error grows unbounded over multi-hour cycles \cite{waag2014critical, rout2024online, dong2026secure}.

While high-fidelity electrochemical physics-based models capture solid-phase diffusion and Butler-Volmer kinetics with high fidelity, their non-linear partial differential algebraic equation (PDAE) structure incurs prohibitive computational complexity \cite{sorouri2025breaking, che2025enhanced, sorouri2026physics}. Solving coupled PDAE systems exceeds the deterministic execution budgets and static memory allocations of standard automotive microcontrollers (e.g., ARM Cortex-M4, Infineon AURIX TC3xx, and NXP S32K) by multiple orders of magnitude \cite{yahyaei2023embedded}.

To overcome physical modeling bottlenecks, recent benchmark studies, most notably the 2025 Stanford relaxation methodology \cite{che2025enhanced}, introduced multi-stage data-driven pipelines utilizing decision tree ensembles (such as sequential 3-stage Random Forests) to capture non-linear relaxation voltage tails during vehicle rest periods. Although tree ensembles improve state tracking across flat OCV intervals, their algorithmic complexity presents severe deployment barriers in production automotive environments \cite{lin2020mcunet, cui2025toward}.

As quantified in Section~\ref{sec:res_g}, the reference Stanford 3-stage Random Forest (RF) baseline \cite{che2025enhanced} requires approximately $1420.0\ \text{kB}$ of Flash memory and contains $\approx 60,000$ parameters. Traversing deep ensemble decision paths induces an inference latency of $1850.0\ \mu\text{s}$ ($1.85\ \text{ms}$) on a standard ARM Cortex-M4 core at 160~MHz. Similarly, sequence-based deep models such as Tabular Transformers ($142.0\ \text{kB}$ Flash, $35.5\text{k}$ parameters) \cite{chen2024transformer} and Vanilla LSTM networks ($53.0\ \text{kB}$ Flash, $13.2\text{k}$ parameters) \cite{xiong2018towards} violate the stringent static memory allocations ($M_{\text{Flash}} \le 16\ \text{kB}$) and worst-case execution time budgets ($T_{\text{exec}} \le 50\ \mu\text{s}$) of commercial BMS electronic control units (ECUs). Furthermore, tree ensembles with irregular branching structures and unconstrained recurrent networks cannot provide static execution determinism, preventing ISO 26262 ASIL-D functional safety certification \cite{yahyaei2023embedded}.

\subsection{Proposed Physics-Informed Knowledge Distillation Paradigm}
\label{sec:intro_c}

To resolve the fundamental trade-off between non-linear representational capacity and embedded execution determinism, Physics-Informed Machine Learning coupled with Knowledge Distillation (PIML-KD) offers an effective paradigm \cite{karniadakis2021physics, li2026condition, amirkhani2021robust, liu2024unsupervised}. Under this architecture, the state estimation problem is decoupled into an offline high-capacity teacher learning stage and an on-chip, physics-constrained student execution engine \cite{amirkhani2021robust, najafi2024statepower, li2026condition}.

As conceptualized in Fig.~\ref{fig:framework_overview}, the teacher model comprises an offline Gradient Boosted Decision Tree (LightGBM) ensemble parameterized with 300 estimators \cite{ke2017lightgbm}. This teacher possesses sufficient non-linear depth to internalize complex, path-dependent electrochemical hysteresis, state-dependent solid-phase diffusion delays, and multi-regime relaxation profiles across the complete operating space without run-time execution limits.

To transfer this non-linear mapping capability into an ultra-efficient edge format, we construct a compact student Multilayer Perceptron (Micro-MLP) featuring a $12 \to 32 \to 16 \to 1$ topology ($961$ total parameters) \cite{lin2020mcunet}. Rather than treating distillation as an unconstrained empirical curve-fitting problem, the student network is strictly regularized via domain-guided physical laws \cite{karniadakis2021physics, sorouri2026physics}: (i) differential relaxation overpotentials ($\Delta V_{10} = V_{\text{mea}}(t) - V_{10}$), (ii) logarithmic diffusion scales ($\tau_{\text{log}} = \ln(t + 1.0)$), (iii) dynamic ohmic dissipation ($P_{\text{ohm}} = R \cdot I_m$), (iv) Arrhenius reaction kinetics ($\Phi_{\text{Arr}} = \exp(-E_a/(R_g T_K))$ with $E_a = 35.0\ \text{kJ/mol}$) \cite{rahimirad2019battery, meng2025domain}, and (v) a composite distillation loss combining Smooth-$L_1$ Huber loss and quadratic physical boundary penalties enforcing $\text{SOC} \in [0.0, 1.0]$. The framework synthesizes directly into a zero-\texttt{malloc} static C header for deterministic execution on automotive electronic control units \cite{yahyaei2023embedded}.

\subsection{Core Contributions and Article Organization}
\label{sec:intro_d}

The principal contributions of this work are summarized as follows:
\begin{enumerate}
	\item \textbf{Physics-Informed Domain-Guided 12-D Feature Engineering:} Formulating a 12-dimensional physical input representation ($X_{\text{PIML}}$) that augments raw sensor telemetry with differential overpotentials ($\Delta V_{10}$), logarithmic diffusion scaling ($\tau_{\text{log}}$), dynamic ohmic dissipation ($P_{\text{ohm}}$), and Arrhenius reaction kinetics ($\Phi_{\text{Arr}}$ with $E_a = 35.0\ \text{kJ/mol}$) \cite{rahimirad2019battery, li2026condition}.
	\item \textbf{Physically Calibrated Hybrid Knowledge Distillation:} Developing a teacher--student compression pipeline wherein a 300-tree LightGBM teacher supervises an ultra-lightweight student Micro-MLP ($12 \to 32 \to 16 \to 1$, $961$ total parameters) under a multi-objective loss enforcing thermodynamic boundary validity \cite{ke2017lightgbm, amirkhani2021robust}.
	\item \textbf{Multi-Regime Generalization and Telemetry Noise Robustness:} Achieving a nominal $25^\circ\text{C}$ MAE of $0.454\%$ and RMSE of $1.236\%$ ($R^2 = 0.9986$). On unseen cells YX07 and YX08, the model achieves MAEs of $2.48\%$ and $3.35\%$, beating the Stanford target ($4.30\%$) \cite{che2025enhanced} by $42.3\%$ and $22.1\%$. The framework enables Coulomb counting resets at $30\ \text{s}$ ($\text{MAE} = 1.170\%$, converging to $0.339\%$ at $300\ \text{s}$) and resists $\pm 5.0\ \text{mV}$ AFE noise ($\text{MAE} = 1.336\%$).
	\item \textbf{MISRA-C Compliant Static Deployment Pipeline:} Providing an automated synthesis pipeline exporting the student Micro-MLP into a static C header (\texttt{bms\_soc\_piml\_mlp3.h}) requiring only $3.75\ \text{kB}$ Flash, $192\ \text{Bytes}$ SRAM, and $12.0\ \mu\text{s}$ latency on an ARM Cortex-M4 at $160\ \text{MHz}$ \cite{yahyaei2023embedded, lin2020mcunet}.
\end{enumerate}

The remainder of this article is structured as follows: Section~\ref{sec:methodology} details the mathematical formulation of the 12 PIML descriptors, teacher pipeline, distillation loss, and static C generation. Section~\ref{sec:results} presents comprehensive experimental results. Section~\ref{sec:conclusion} concludes the article.